\documentclass{beamer}
\usepackage[T1]{fontenc}
\usepackage[utf8]{inputenc}
\usepackage[round]{natbib}
\usepackage{graphicx}
\usetheme{CambridgeUS}  % You can change theme as needed

\title{Bitter Pills to Swallow:\\ The Enforcement Costs of Health Litigation}
\author{Darcio Genicolo-Martins \and Paulo Furquim de Azevedo}
\institute{Insper Institute of Education and Research,\\ Brazil}
\date{}  % No date specified

\begin{document}

\begin{frame}
  \titlepage
\end{frame}

\section{Introduction}

\begin{frame}{Introduction}
  \begin{itemize}
    \item The right to health is constitutional in Brazil, \\leading many patients to sue the government for medications \citep{wang2015right}.
    \item Health litigation cases have increased significantly, \\forcing unplanned spending outside the regular health budget \citep{cnj2019judicializacao}.
  \end{itemize}
\end{frame}

\begin{frame}{Motivation}
  \begin{itemize}
    \item Court-mandated urgent procurement bypasses standard tender procedures.
    \item Urgent purchases often lack competitive bidding, \\potentially raising prices.
    \item Emergency procurement and weak oversight can increase inefficiency \\and corruption risks \citep{basheka2011economic, mironov2016corruption}.
  \end{itemize}
\end{frame}

\section{Research Questions}

\begin{frame}{Research Questions}
  \begin{itemize}
    \item What are the enforcement costs (budgetary and efficiency costs) \\of court-mandated health procurement?
    \item How do prices paid for medicines differ between litigated purchases \\and normal procurement?
    \item What are the policy implications for balancing health rights \\and fiscal sustainability?
  \end{itemize}
\end{frame}

\section{Context and Data}

\begin{frame}{Data \& Methodology}
  \begin{itemize}
    \item \textbf{Data:} Drug procurement records from S\~ao Paulo state’s health department, \\including identification of court-mandated purchases.
    \item \textbf{Method:} Compare prices and other outcomes for the same medicines \\across procurement types (ordinary vs. litigated).
    \item \textbf{Approach:} Regression analysis controlling for drug fixed effects \\to isolate the price difference attributable to litigation enforcement.
  \end{itemize}
\end{frame}

\section{Theory}

\begin{frame}{Theoretical Predictions}
  \begin{itemize}
    \item Ordinary: Ample competition, price $\approx c$ (marginal cost).
    \item Urgent: Reduced competition, price increases to $c + \delta$ (markup $\delta$).
    \item Litigated: Enforcement costs, price rises to $c + \delta + \pi\gamma$ \\(penalty cost term $\pi\gamma$).
  \end{itemize}
\end{frame}

\section{Results}

\begin{frame}{Rising Health Litigation in Brazil}
  \begin{center}
    \includegraphics[width=\textwidth]{figures/figure1.png}
  \end{center}
  \footnotesize Health-related lawsuits have increased dramatically in recent years \citep{cnj2019judicializacao}.
\end{frame}

\begin{frame}{Budgetary Impact of Litigation}
  \begin{center}
    \includegraphics[width=\textwidth]{figures/figure2.png}
  \end{center}
  \footnotesize Litigated purchases constitute a growing portion of total \\pharmaceutical spending in the public health sector.
\end{frame}

\begin{frame}{Comparison of Procurement Channels}
  \begin{center}
    \includegraphics[width=\textwidth]{figures/figure3.pdf}
  \end{center}
  \scriptsize \textit{Note:} "Administrative" refers to urgent procurement without a court order (internal administrative process).
\end{frame}

\begin{frame}{Result: Higher Prices in Litigated Procurement}
  \begin{center}
    \includegraphics[width=\textwidth]{figures/figure4.png}
  \end{center}
  \footnotesize Medicines procured via litigation cost significantly more on average (price premium).
\end{frame}

\begin{frame}{Result: Smaller Purchase Quantities}
  \begin{center}
    \includegraphics[width=\textwidth]{figures/figure5.png}
  \end{center}
  \footnotesize Litigated purchases typically involve much smaller quantities per order than regular procurement.
\end{frame}

\begin{frame}{Result: Faster Delivery Times}
  \begin{center}
    \includegraphics[width=\textwidth]{figures/figure6.png}
  \end{center}
  \footnotesize Litigated purchases are delivered much faster (often within days) compared to ordinary procurement (which can take months).
\end{frame}

\begin{frame}{Result: Low Supplier Competition}
  \begin{center}
    \includegraphics[width=\textwidth]{figures/figure7.png}
  \end{center}
  \footnotesize Litigated purchases are often single-sourced, indicating very low competition among suppliers.
\end{frame}

\section{Conclusion}

\begin{frame}{Conclusion}
  \begin{itemize}
    \item Enforcing the right to health via litigation ensures patient access \\but incurs high procurement costs.
    \item Litigated purchases show higher prices and inefficiencies \\(smaller batches, single suppliers).
    \item Policymakers should improve regular health provision to reduce reliance \\on lawsuits and associated costs.
  \end{itemize}
\end{frame}

\section*{References}

\begin{frame}[allowframebreaks]{References}
  \footnotesize
  \bibliographystyle{apalike}
  \bibliography{references}
\end{frame}

\end{document}
